\documentclass[CJKmath,a4paper,10pt]{ctexart}
\usepackage{geometry}
\geometry{inner=1cm,outer=1.5cm,top=2.45cm,bottom=1cm}
\usepackage[dvipsnames,svgnames,x11names,table]{xcolor}


\RequirePackage{amssymb} % Must be loaded before unicode-math
\RequirePackage{unicode-math} % Math fonts in xetexorluatex
%\setmathfont{texgyrepagella-math.otf}
%\setmainfont{STIXTwoText}[
%	Path=./fonts/STIXTwoText/,
%  Extension = .otf,
%  UprightFont = STIXTwoMath-Regular,
%  BoldFont = STIXTwoText-SemiBold,
%  ItalicFont = STIXTwoText-Italic,
%  BoldItalicFont = STIXTwoText-BoldItalic,
%]
%\setmathfont{STIXTwoMath-Regular.otf}

\usepackage{tikz,calc}
\usetikzlibrary{calc,arrows,shadows.blur,tikzmark,patterns,decorations.pathmorphing,arrows.meta}
% By default all math in TikZ nodes are set in inline mode. Change this to
% displaystyle so that we don't get small fractions.
\everymath{\displaystyle}

\usepackage[most]{tcolorbox}
\tcbuselibrary{breakable, skins,theorems}


\usepackage{xkeyval}
\makeatletter 
\usepackage{amsmath,amssymb,amsthm}
\usepackage{xkeyval}
\tcbset{
	common/.style={
		fontupper=\rmfamily,
		lower separated=true,
		coltitle=black,
		colback=gray!5,
		boxrule=0.5pt,
		fonttitle=\bfseries,
		enhanced,
		breakable,
		top=8pt,
		before skip=8pt,
		attach boxed title to top left={yshift=-0.05in,xshift=0.15in},
		boxed title style={
			boxrule=0pt,
			colframe=black,
			arc=0pt,
			outer arc=0pt,
			drop fuzzy shadow
		},
		separator sign={.},
		drop fuzzy shadow
	},
	litistyle/.style={  %例题风格
		common,
		colframe=gray,% 
		colback=white,
		colbacktitle=Gold,
		sharp corners,rounded corners=southeast,
		overlay unbroken and last={
			\node[anchor=south east, outer sep=0pt] at (\linewidth-width,0){\textcolor{black!60}{$\blacksquare$}};}
	},
	% -----------------------------------------------------------------------------
	% Explanation: BreakableBox@title key
	% The following defines a pgfkeys/tcolorbox key named
	%   "BreakableBox@title" which accepts 2 arguments (see 
	%   "/.code n args={2}"). When the key is used like
	%   "BreakableBox@title={<env>}{<optTitle>}" the code below is executed
	%   with #1=<env> and #2=<optTitle>.
	%
	% Purpose:
	% - If the second argument (#2) is empty, set the tcolorbox title to
	%     \csname <env>name\endcsname~\thetcbcounter
	%   which expands to the environment name (e.g. "liti" -> \litiname)
	%   followed by the current counter value.
	% - If the second argument is provided, append " (#2)" to the title.
	%
	% Notes:
	% - This relies on \makeatletter being active so that keys with '@'
	%   are allowed in control sequence names.
	% - \ifblank is used to test whether #2 is empty; ensure the package
	%   that provides \ifblank (for example, etoolbox or xparse) is loaded
	%   earlier if necessary.
	% - Example usage inside the code: \tcbset{title={...}} sets the
	%   tcolorbox title key accordingly.
	%
	BreakableBox@title/.code n args={2}
	{
		\ifblank{#2}
		{\tcbset{title={\csname #1name\endcsname~\thetcbcounter}}}
		{\tcbset{title={\csname #1name\endcsname~\thetcbcounter\ (#2)}}}
	},
}      
% define an internal control sequence \newbreakablebox for fancy mode's newtheorem
% #1 is the environment name, #2 is the prefix of label, #3 is the style
% style: thmstyle, defstyle, prostyle
% e.g. \newbreakablebox{theorem}{thm}{thmstyle}
% will define two environments: numbered ``theorem'' and no-numbered ``theorem*''
% WARNING FOR MULTILINGUAL: this cs will automatically find \theoremname's definition,
% WARNING FOR MULTILINGUAL: it should be defined in language settings.
\newcommand{\newbreakablebox}[3]{
	\ifcsundef{#1name}{%
		% define a default name if not defined before
		\tcbset{BreakableBox@title/.code n args={2}{\tcbset{title={use newcommand define #1name}}} }
	}{\relax}
	\DeclareTColorBox[auto counter,number within=section]{#1}{ g o t\label g }{
		common, % use common style
		#3, % use the style passed in
		IfValueTF={##1}
		{BreakableBox@title={#1}{##1}}
		{
			IfValueTF={##2}
			{BreakableBox@title={#1}{##2}}
			{BreakableBox@title={#1}{}}
		},
		IfValueT={##4}
		{
			IfBooleanTF={##3}
			{label={##4}}
			{VIVID@label={#2}{##4}}
		}
	}
	\DeclareTColorBox{#1*}{ g o }{
		common,#3,
		IfValueTF={##1}
		{BreakableBox@title={#1}{##1}}
		{
			IfValueTF={##2}
			{BreakableBox@title={#1}{##2}}
			{BreakableBox@title={#1}{}}
		},
	}
}
% define several environment 
% we define headers like \definitionname before
\newcommand{\litiname}{例题}
\newbreakablebox{liti}{def}{litistyle}


%补充内容
%%设置新字体
%%定义带圈数字命令
\newfontfamily{\nmfont}{circlenumber}
[%
Extension=.otf,
Path=./fonts/]

\newcommand{\quan}[1]{{\nmfont \symbol{#1}}}
\newcommand{\kk}[1]{\quan{\numexpr32+#1}}%\kk{<参数范围1-95>}96、97、98、99分别用\quan{196} \quan{197} \quan{199} \quan{201}

%脚注使用带圈数字
\newcommand*\kkctr[1]{%
	\protect\kk{\number\numexpr\value{#1}\relax}}
\renewcommand*\thefootnote{\textcolor{black}{\kkctr{footnote}}}

%%无悬挂脚注格式
\renewcommand\@makefntext[1]{%
	\setlength\parindent{2\ccwd}\selectfont
	\@thefnmark\ #1}

%修改\part，使其不分页
\def\@endpart{%
	\thispagestyle{empty}
	\vskip40\p@%
	\@afterheading}



\RequirePackage{enumitem}
%\newenvironment{myenum}{\begin{enumerate}[label=\protect\kk{\arabic*}]\small}{\end{enumerate}}%
\setlist{noitemsep}
\setlist[enumerate, 1]{label=\protect\kk{\arabic*},itemsep=0.5ex}  
\makeatother



%%%marker环境
\newtcolorbox{marker}[1][]{enhanced,before skip=2mm,
	after skip=3mm,fontupper=\rmfamily,
	boxrule=0.4pt,left=5mm,right=2mm,top=1mm,bottom=1mm,
	colback=yellow!50,colframe=yellow!20!black,
	sharp corners,rounded corners=southeast,
	arc is angular,arc=3mm,underlay={%
		\path[fill=tcbcolback!80!black] ([yshift=3mm]interior.south east)--++(-0.4,-0.1)--++(0.1,-0.2);
		\path[draw=tcbcolframe,shorten <=-0.05mm,shorten >=-0.05mm] ([yshift=3mm]interior.south east)--++(-0.4,-0.1)--++(0.1,-0.2);
		\path[fill=yellow!50!black,draw=none] (interior.south west) rectangle node[white]{\Huge\bfseries !} ([xshift=4mm]interior.north west);
	},
	drop fuzzy shadow,#1
}


\usepackage{graphicx}
\usepackage{float}
\begin{document}
	\section{物理习题}
\begin{liti}	
 如图甲所示，固定粗糙斜面的倾角为$37^\circ$，与斜面平行的轻弹簧下端固定在$C$处，上端连接质量为$1\mathrm{kg}$的小滑块（视为质点），$BC$为弹簧的原长。现将滑块从$A$处由静止释放，在滑块从释放至第一次到达最低点的过程中，其加速度$a$随弹簧的形变量$x$的变化规律如图乙所示（取沿斜面向下为加速度的正方向），取$\sin37^\circ=0.6$，$\cos37^\circ=0.8$，$g=10\mathrm{m/s^2}$。根据上述过程，试求出下列各量：
 \begin{enumerate}
	\item 滑块与斜面间的动摩擦因数$\mu$；
	\item 当$x=0$时滑块的动能$E_\text{k}$；
	\item 弹簧的最大弹性势能$E_\text{p}$。
\end{enumerate}


% 题目配图（图甲：斜面+弹簧+滑块）
\begin{figure}[H]
\centering
\begin{tikzpicture}[>=Stealth,scale=1]
	% 斜面与地面
	\begin{scope}[xshift=-5cm]
	\coordinate (O) at (0,0);
	\coordinate (D) at (4,0);
	\coordinate (S) at (4,3);
	\draw[thick] (O) -- (D);
	\draw[thick] (O) -- (S);

	% 倾角标注
	\draw (1.2,0) arc[start angle=0,end angle=37,radius=1.2];
	\node at (1.55,0.45) {$37^\circ$};

	% 斜面上的关键点
	\coordinate (C) at (0,0);
	\coordinate (B) at (3,2.25);
	\coordinate (A) at (4,3);

	% 弹簧与滑块
	\draw[decorate,decoration={coil,segment length=4pt,amplitude=2.4pt}] ($(C)+(0,0.2)$) -- ($(A)+(0,0.2)$);
	\begin{scope}[shift={(A)},rotate=37]
		\draw[fill=white] (-0.28,0) rectangle (0.28,0.42);
	\end{scope}

	% 点与文字标注
	\fill (A) circle (1.2pt);
	\fill (B) circle (1.2pt);
	\fill (C) circle (1.2pt);
	\node[above left,rotate=37] at ($(A)+(0.2,0.1)$) {$A$};
	\node[above,rotate=37] at ($(B)+(-0.1,0.3)$) {$B$};
	\node[below left] at (C) {$C$};
	\node[right] at ($(A)+(0.35,0.1)$) {滑块};

	% BC 为原长
	\draw[|<->|] ($(B)+(-0.32,0.43)$) -- ($(C)+(-0.32,0.43)$)
		node[midway,sloped,above] {$BC$为原长};

	\node at (2.4,-1.5) {图甲-斜面装置};
\end{scope}
	\begin{scope}[x=10cm,y=1cm,scale=0.5,xshift=10cm,yshift=2cm]
	% 图乙：a-x图像
	
	%\begin{tikzpicture}
		\coordinate (A) at (0.2,4);
		\coordinate (B) at (0,2);
		\coordinate (C) at (0.2,0);
		\coordinate (D) at (0.6,-4);
		\coordinate (E) at(0,4);
		\coordinate (F) at(0,2);
		% 坐标轴
		\draw[->] (0,-4.5) -- (0,4.5) node[above] {$a\ (\mathrm{m\cdot s^{-2}})$};
		\draw[->] (-0.1,0) -- (0.7,0) node[right] {$x\ (\mathrm{m})$};
		% 刻度
		\foreach \y in {-4,-3,-2,-1,1,2,3,4} {
			\draw (0.02,\y) -- (-0.02,\y) node[left] {$\y$};
		}
		\foreach \x in {0.2,0.4,0.6} {
			\draw (\x,0.05) -- (\x,-0.05) node[below] {$\x$};
		}
		\draw(A)--(B)--(C)--(D);
		\draw[dotted,red] (A) node[right]{A}--(E);
		\node[right] at(F){B};
		\node[above]at(0.2,0){Q};
		\node[below left]at(0,0){O};
		\node at (0.3,-5) {图乙-$a-x$变化规律};
\end{scope}		
\end{tikzpicture}
	\caption{\emph{图甲：斜面装置；图乙：$a-x$变化规律}}
\end{figure}
\tcblower
\begin{enumerate}
	\item \textbf{求滑块与斜面间的动摩擦因数$\mu$}\\
	当$x=0$时，弹簧处于原长，弹力为$0$。对滑块沿斜面方向由牛顿第二定律：$mg\sin37^\circ - \mu mg\cos37^\circ = ma_0$
代入$m=1\mathrm{kg}$，$g=10\mathrm{m/s^2}$，$\sin37^\circ=0.6$，$\cos37^\circ=0.8$，$a_0=2\mathrm{m/s^2}$：
\[1\times10\times0.6 - \mu\times1\times10\times0.8 = 1\times2 \Rightarrow \mu = 0.5\]


\item \textbf{当$x=0$时滑块的动能$E_\text{k}$}

根据\textbf{动能定理}，合外力做功等于动能变化：$W_\text{合} = \Delta E_\text{k}$。
合外力$F_\text{合}=ma$，在$a-x$图像中，合外力做功等于$m$乘以$a-x$图线与$x$轴围成的面积（高中面积法）。

从$A$（$x=0.2\mathrm{m}$）到$B$（$x=0$），$a-x$图为梯形：
上底$a_1=2\mathrm{m/s^2}$，下底$a_2=4\mathrm{m/s^2}$，高$\Delta x=0.2\mathrm{m}$
\[
S = \frac{(a_1+a_2)}{2} \cdot \Delta x = \frac{(2+4)}{2} \times 0.2 = 0.6\ \mathrm{m\cdot s^{-2}}
\]
合外力做功：
$
W_\text{合} = mS = 1\times0.6 = 0.6\ \mathrm{J}
$
因滑块从静止释放，初动能为$0$，故$E_\text{k} = W_\text{合} = 0.6\mathrm{J}$。

\item \textbf{求弹簧的最大弹性势能$E_\text{p}$}

当滑块第一次到达最低点时，速度为$0$，弹簧形变量最大，设为$x_m$。
由图乙中$x>0$部分知，滑块压缩弹簧时的加速度满足
\[
a=2-10x.
\]
\textbf{思路：} 
	\emph{由于在全部过程中，动能变化量为$\Delta E = 0$（初末速均为0），由第二问知：物体由A到B(弹簧伸长量为0)点动能等于梯形ABOQ的面积为0.6，所以从B点到弹簧压缩量最大点$X_m$与x轴围成的面积（射线BQ与x轴围成的面积）应为-0.6。
	}
	
从$x=0$到第一次到达最低点，由动能定理可得合外力做功等于动能变化。
在$a-x$图像中，这对应于图线与$x$轴围成的有向面积乘以质量。

其中，$x=0$到$x=0.2$之间图线在$x$轴上方，围成一个三角形，面积为
\[
S_1=\frac12\times 0.2\times 2=0.2.
\]

$x=0.2$到$x=x_m$之间图线在$x$轴下方，围成一个三角形，面积为
\[
S_2=\frac12\times (x_m-0.2)\times (10x_m-2).
\]

因为此过程末动能为$0$，初动能为$E_\text{k}=0.6\mathrm{J}$，所以合外力做功为
\[
W_\text{合}=0-E_\text{k}=-0.6\mathrm{J}.
\]

又$m=1\mathrm{kg}$，故有
\[
S_1-S_2=-0.6.
\]

即
\[
0.2-\frac12(x_m-0.2)(10x_m-2)=-0.6.
\]

化简得
\[
5x_m^2-2x_m-0.6=0.
\]

整理得
\[
x_m=0.6\mathrm{m}
\]
（负值舍去）。

故弹簧的最大弹性势能为
\[
E_\text{p}=\frac12 kx_m^2=\frac12\times10\times(0.6)^2=1.8\mathrm{J}.
\]
\end{enumerate}
\noindent
\textbf{最终答案：}
(1) $\boxed{\mu=0.5}$；(2) $\boxed{E_\text{k}=0.6\mathrm{J}}$；(3) $\boxed{E_\text{p}=1.8\mathrm{J}}$

\end{liti}
\end{document}